% =====================================================================
%  FlexTrack-V2 — Introduction (IEEE TPAMI / IEEEtran)
%  Journal extension of "What You Have is What You Track" (ICCV 2025)
% =====================================================================

\section{Introduction}
\label{sec:intro}

\IEEEPARstart{M}{ultimodal} visual object tracking augments an RGB camera with a complementary sensor (thermal, depth, or event) to ensure reliability when a single modality fails under poor illumination, cluttered backgrounds, or fast motion. While recent trackers have steadily pushed accuracy upward on dedicated complete-modality benchmarks, they rely on a tacit assumption that rarely holds in deployment: that both modalities are fully and synchronously available at every frame. In real-world sensor rigs, this assumption is constantly violated. Cameras de-synchronize, depth returns drop out on distant or reflective surfaces, and thermal feeds stall, rendering the auxiliary modality intermittently missing within a video.

When this intermittent missingness occurs, conventional multimodal trackers degrade sharply. Their fusion modules are trained on complete modality pairs and exhibit no principled behaviour when one stream vanishes. Feeding zeros or repeating the last frame injects noise, while rigid template-matching pipelines cannot re-allocate computation to compensate. Prior attempts to improve robustness treat missing modalities as a set of special cases, using separate prompts or branches for each degradation, which neither scales to the combinatorial space of dropout patterns nor adapts the fusion complexity to the degree of missingness.

Our conference work, FlexTrack~\cite{tan2025flextrack}, addressed this missing-modality challenge by adapting its complexity to modality availability based on the principle of \emph{what you have is what you track}. Instead of a fixed fusion cost, FlexTrack routes each video clip through a Heterogeneous Mixture-of-Experts (HMoE) whose experts span a range of capacities, driven by the current data. Trained with a video-level masking strategy to preserve temporal consistency, it yielded a single unified model that handles all three modality pairs and remains robust as the missing rate grows.

Yet, adapting capacity to absence is only a partial solution. While allocating fewer experts avoids wasting computation when a modality is missing, it does nothing to recover the information that the absent sensor would have provided. Because the two modalities of a scene are strongly correlated, the present stream in fact carries a strong prior about the absent one. FlexTrack, like earlier robust trackers, fundamentally \emph{copes with} an absent stream rather than actively \emph{compensating for} it, leaving this cross-modal prior unused.

In this paper, we propose FlexTrack-V2, which extends our framework from \emph{adaptive coping} to \emph{active compensation}. Our first innovation is a Bilateral Modality-specific feature Reconstruction and Hallucination (BMR) module. Rather than masking an absent modality, two lightweight cross-modal projectors hallucinate its features from the available stream, with supervision from a self-supervised reconstruction loss whenever the target modality is genuinely observed. Coupled with the HMoE, this BMR-HMoE fusion both sizes its computation to the input and actively fills in the missing gaps. Our second innovation is Curriculum Missing Augmentation (CMA) with online self-distillation, which stabilizes the training of such a model. By annealing the missing probability over training and distilling a complete-modality teacher into a degraded student of the same network, CMA ensures that reconstructed predictions stay faithful to complete-modality ones. These components are instantiated on a unified architecture comprising a shared Fast-iTPN encoder and a Mamba (selective state-space) interaction neck, trained jointly across RGB-T, RGB-D, and RGB-E data.

Extensive experiments across nine benchmarks demonstrate the effectiveness of FlexTrack-V2. Under complete modalities, the model attains parity with the state-of-the-art conference version. Under missing modalities, the regime that matters most for deployment, FlexTrack-V2 improves consistently and substantially. The largest gains appear precisely on the benchmarks and dropout settings where active reconstruction and cross-completeness consistency are expected to help, exhibiting the ideal behaviour of a robust tracker: no cost when everything is present, and real recovery when something is not.

\medskip
\noindent Contributions. This work makes the following contributions:
\begin{itemize}
  \item We recast robustness to missing modalities from \emph{adaptive coping} to \emph{active compensation} and propose BMR-HMoE, a fusion module that unifies capacity-adaptive heterogeneous experts with bilateral cross-modal feature hallucination.
  \item We propose Curriculum Missing Augmentation with online self-distillation (CMA), which anneals the missing rate over training and aligns representations across completeness levels to stabilize the learning of the reconstruction pathway.
  \item We instantiate these ideas in a single unified tracker (using a single set of parameters for RGB-Thermal, RGB-Depth, and RGB-Event tracking under both complete and missing protocols) built on a Fast-iTPN encoder with an Mamba interaction neck.
  \item We conduct a comprehensive study across complete and missing settings, demonstrating that FlexTrack-V2 matches the prior state of the art on complete modalities and improves markedly on missing ones.
\end{itemize}

\medskip
\noindent Relation to the conference version. A preliminary version of this work appeared at ICCV~2025~\cite{tan2025you}. That paper introduced the adaptive-complexity principle, the heterogeneous MoE fusion, and the video-level masking strategy. The present paper substantially extends it along four axes: (i) it adds the BMR module, moving from passive masking to active bilateral reconstruction of missing features; (ii) it introduces the curriculum missing-augmentation and online self-distillation training scheme; (iii) it consolidates the method on a unified Fast-iTPN\,+\,Mamba architecture; and (iv) it provides an expanded experimental study with new benchmarks, additional missing-modality protocols, and thorough ablations. Together these changes constitute more than a $30\%$ increase in new technical content and lead to consistent improvements over the conference model in the missing-modality regime.
